\documentclass[10pt,landscape,a4paper]{article}
\usepackage[ngerman]{babel}
\usepackage[T1]{fontenc}
\usepackage[dvipsnames]{xcolor}
%\usepackage[LY1,T1]{fontenc}
\usepackage{tikz}
\usetikzlibrary{shapes,positioning,arrows,fit,calc,graphs,graphs.standard}
\usepackage[nosf]{kpfonts}
\usepackage{multicol}
\usepackage{wrapfig}
\usepackage[top=2mm,bottom=3mm,left=2mm,right=2mm,headsep=1mm,includehead]{geometry}
\usepackage[framemethod=tikz]{mdframed}
\usepackage{microtype}
\usepackage{pdfpages}
\usepackage{csquotes}
\usepackage{booktabs}

\usepackage{amsmath}
\usepackage{bbm}
\newcommand{\N}{\mathbb{N}}
\newcommand{\Z}{\mathbb{Z}}
\newcommand{\R}{\mathbb{R}}
\newcommand{\C}{\mathbb{C}}
\newcommand{\F}{\mathbb{F}}
\newcommand{\K}{\mathbb{K}}
\renewcommand{\S}{\mathbb{S}}
\DeclareMathOperator{\im}{Im}
\DeclareMathOperator{\Ker}{Ker}
\DeclareMathOperator{\Sym}{Sym}
\DeclareMathOperator{\sgn}{sgn}
\DeclareMathOperator{\Hom}{Hom}
\DeclareMathOperator{\lc}{lc}
\DeclareMathOperator{\Mat}{Mat}
\DeclareMathOperator{\id}{id}
\DeclareMathOperator{\ggT}{ggT}
\DeclareMathOperator{\kgV}{kgV}

\newcommand\mat[1]{
	\begin{pmatrix}
		#1
	\end{pmatrix}
}

\usepackage{blindtext}

\usepackage{titlesec}
\input{../layout.tex}

\usepackage{hyperref}
\newcommand\nohyper[1]{\begin{NoHyper}#1\end{NoHyper}}

\usepackage{lastpage,fancyhdr}
\pagestyle{fancy}
\fancyhf{}
\fancyhead[L]{Rasmus Buurman}
\fancyhead[C]{\textsc{Algebraische Strukturen}}
\fancyhead[R]{Seite \thepage{}/\nohyper{\pageref{LastPage}}}
\renewcommand\headrulewidth{0.5pt}

\let\bar\overline

\begin{document}
\small
\begin{multicols*}{4}
	\input{contents/gruppen-homomorphismen.tex}
	\section*{Faktorgruppen}

\subsection*{Linksnebenklassen}
Gruppe $G$ und $U\le G$. $g,h\in G$.

$ g \sim h \iff g^{-1}h\in U,\quad \bar g = gU$\\
$ G/U = \{gU \mid g \in G \}, \quad |G:U| = |G/U| $

\subsection*{Satz von Lagrange}
Endl. Gruppe $G$ und $U \le G$.

$|G| = |U|\cdot|G:U|$

\subsubsection*{Produktformel}
$U,V$ endlich.
$|U \cdot V| = \frac{|U| \cdot |V|}{|U \cap V|}$

\subsection*{Ordnung eines Elements}
$G$ Gruppe, $g\in G$.

Ordnung $o(g):=|\langle g \rangle|$ teilt $|G|$.

Außerdem: $o(g)=\inf\{k>0 \mid g^k = e\}$

\subsection*{Normalteiler}
$U \le G$ Normalteiler: $\forall g\in G, u\in U: gug^{-1}\in U$.

\subsubsection*{Abelsche Gruppen und Normalteiler}
$G$ abelsch $\implies U \le G$ normal

\subsubsection*{Äquivalente Kriterien}
\begin{itemize}
	\item $\forall g\in G: gUg^{-1} = U$
	\item $\forall g \in G: gU = Ug$
	\item $\forall g,h\in G: (gU) \cdot (hU) = ghU$
\end{itemize}

\subsubsection*{Kern ist normal}
$\alpha: G \to H \in \Hom \implies \Ker(\alpha) \trianglelefteq G$

\subsubsection*{Weitere Eigenschaften}
$U \le G$ und $N \trianglelefteq G$. Dann:

$UN \le G$, $N \trianglelefteq UN$, $U \cap N \trianglelefteq U$

\subsection*{Faktorgruppe $G/U$}
$U \trianglelefteq G$. Dann ist $G/U$ eine Gruppe.

	\section*{Homomorphiesatz}
$\alpha: G \to H \in \Hom$.

$\tilde{\alpha}: G/\Ker\alpha \to \im\alpha: \bar g \mapsto \alpha(g)$ Isomorphismus.

$\implies G/\Ker\alpha \cong \im\alpha$

\subsection*{1. Isomorphiesatz}
$U \le G$, $N \trianglelefteq G$. Dann: $U/U \cap N \cong UN/N$.

\subsection*{2. Isomorphiesatz}
$M \subseteq N \subseteq G$, beide normal. Dann $N/M \trianglelefteq G/M$ und:

$(G/M)/(N/M) \cong G/N$

	\section*{Zyklische Gruppen}

\subsection*{Isomorphismus in zyklischen Gruppen}
$G = \langle g \rangle$.

$|G|=\infty \implies \alpha: \Z \overset{\cong}{\to} G: z \mapsto g^z$

$|G| = n \implies \bar{\alpha} : \Z_n \overset{\cong}{\to} G: \bar{z} \mapsto g^z$


	\input{contents/ringe-koerper.tex}
	\section*{Teilbarkeit}

\subsection*{Nullteiler}
$\exists 0 \neq b\in R: a\cdot b = 0$
($R$ kommutativer Ring)

\subsection*{Integritätsbereich}
$0$ ist einziger Nullteiler.

\begin{itemize}
	\item Körper sind Integritätsbereiche
	\item Unterringe von Integritätsbereichen sind Integritätsbereiche
	\item $R$ Integritätsbereich $\implies R[t]$ auch, $R[t]^* = R^*$
\end{itemize}

In Integritätsbereichen kann man kürzen.

\subsection*{$\ggT$ und $\kgV$}

\subsubsection*{$\ggT$}
\begin{enumerate}
	\item $g | a$ und $g | b$
	\item $\forall h \in R$ mit $h | a$ und $h | b$ gilt $h | g$
\end{enumerate}

\subsubsection*{$\kgV$}
\begin{enumerate}
	\item $a | k$ und $b | k$
	\item $\forall l \in R$ mit $a | l$ und $b | l$ gilt $k | l$
\end{enumerate}

\subsubsection*{Fast eindeutig}
$\ggT$ und $\kgV$ sind bis auf Multiplikation mit Einheiten eindeutig bestimmt.

\subsection*{Teilbar}
$b \mid a \iff \langle a \rangle \subseteq \langle b \rangle$

	\section*{Faktorieller Ring}
Jedes $0 \neq a \in R \setminus R^*$ lässt sich als Produkt von endlich vielen Primelementen schreiben.

\subsection*{Irreduzibel und prim}
$0 \neq p \in R \setminus R^*$.

$p$ irreduzibel: $p = a \cdot b \implies a \in R^* \lor b \in R^*$

$p$ prim: $p \mid a \cdot b \implies p \mid a \lor p \mid b$

\subsubsection*{Prim $\implies$ irreduzibel}
(In faktoriellen Ringen auch andersrum)

	\section*{Euklidische Ringe}
Es gibt $\nu: R \setminus \{0\} \to \N$, so dass $a = q \cdot b + r$ mit entweder $r = 0$ oder $\nu(r) < \nu(b)$.

\begin{itemize}
	\item $\Z$ ist euklidisch mit $|\cdot|$
	\item $K[t]$ ist euklidisch mit $\deg$ ($K$ ist Körper)
\end{itemize}

	\section*{Hauptidealringe}
Integritätsbereich $R$ ist Hauptidealring, wenn jedes Ideal von einem Element erzeugt wird.

\begin{itemize}
	\item Euklidisch $\implies$ Hauptidealring
	\item $K$ Körper $\implies$ $K[t]$ Hauptidealring
	\item Hauptidealring $\implies$ faktoriell
\end{itemize}

	\section*{Euklidischer Algorithmus}

\subsection*{Beispiel: Euklidischer Algorithmus}
$a=48$ und $b=-30$
\begin{align*}
	48 &= (-1) \cdot (-30) + 18 \\
	-30 &= (-2) \cdot 18 + 6 \\
	18 &= 3 \cdot 6 + 0
\end{align*}
Als letzter Rest $\neq 0$ ist $\ggT(48,-30) = 6$.

\subsection*{Erweiterter euklidischer Algorithmus}

\subsubsection*{Lemma von B\'ezout}
$a,b\in\Z$, nicht beide 0.

Dann $\exists s,t\in\Z: \ggT(a,b)=sa + tb$.

\subsubsection*{Manuelle Methode}
\begin{enumerate}
	\item Berechne $\ggT(a,b)$
	\item Schrittw. rückw. einsetzen
\end{enumerate}

\subsection*{Anwendung}

\subsubsection*{Multiplikatives Inverse in $\Z_n$}
$0\neq z \in \Z_n$ hat Inv. $\iff \ggT(z,n)=1$

EEA liefert Inv. als $s$.

\subsubsection*{Lineare diophantische Gleichung}
\[c_0 = c_1\cdot x_1 + \ldots + c_n \cdot x_n\]
hat Lösung $\iff \ggT(c_1,\ldots,c_n) \mid c_0$.

	\section*{Chinesischer Restsatz}
$m_i\in\N$ pw. teilerfremd, $a_i\in\Z$.

Dann $\exists x, 0 \le x < M := \Pi m_i$ mit:

$x \equiv a_1 \pmod{m_1}, \ldots, x \equiv a_n \pmod{m_n}$

Lösung eindeutig modulo $M$.

\subsection*{Berechnung}
$M_i := \frac{M}{m_i}$. EEA für: $t_i m_i + s_i M_i = 1$.

Dann: $x = \big(\sum a_i s_i M_i\big) \mod M$.

\subsection*{Modulos zusammenfassen}
$m_1,m_2\in\N$ teilerfremd.

$x \equiv 1 \pmod{m_1} \land x \equiv 1 \pmod{m_2}$

$\implies x \equiv 1 \pmod{m_1\cdot m_2}$

	\section*{Tipps \& Tricks}

\subsection*{Ordnung}
$o(g^m) = \frac{\kgV(m,o(g))}{|m|}$

\subsubsection*{In $\Z_n$}
$\bar m \in \Z_n$ Dann: $o(\bar m) = \frac{\kgV(m,n)}{m}$

\subsection*{Polynome}
\subsubsection*{Nullpolynom}
Hat ein Polynom mehr paarweise verschiedene Nullstellen als sein Grad,
ist es das Nullpolynom.

\subsection*{Permutationen}

$\displaystyle \sigma = \mat{1 & 2 & 3 & 4 & 5 \\ 3 & 4 & 5 & 2 & 1}$

\subsubsection*{Zykelschreibweise}
$\sigma = (1~3~5) \circ (2~4)$

\subsubsection*{Transpositionsschreibweise}
$\sigma = (1~3) \circ (3~5) \circ (2~4)$

Benachbarte Zahlen: \\
$\sigma = (1~2) \circ (2~3) \circ (2~1) \circ (3~4) \circ (4~5) \circ (4~3) \circ (2~3) \circ (3~4) \circ (3~2)$

\subsection*{Ringe Klassifikation}
$R$ Körper $\iff R[t]$ eukl.$\iff R[t]$ Hauptidealring

\subsection*{Einheiten in $\Z_n$}
$m$ ist Einheit, wenn $\ggT(m,n) = 1$.

\end{multicols*}
\end{document}
